\documentclass[UTF8,11pt,a4paper]{ctexart}
\usepackage[left=1.6cm,right=1.6cm,top=2cm,bottom=2cm]{geometry}
\usepackage{amsmath,amssymb,graphicx,longtable,booktabs,array,calc,hyperref,xurl,setspace,fancyvrb,needspace}
\setmainfont{Times New Roman}
\setmonofont{Menlo}
\setCJKmainfont{Songti SC}
\setCJKmonofont{Songti SC}
\hypersetup{hidelinks}
\setlength{\LTleft}{0pt plus 1fill}
\setlength{\LTright}{0pt plus 1fill}
\setlength{\emergencystretch}{3em}
\setlength{\tabcolsep}{4pt}
\setstretch{1.12}
\AddToHook{cmd/section/before}{\Needspace{7\baselineskip}}
\providecommand{\tightlist}{\setlength{\itemsep}{0pt}\setlength{\parskip}{0pt}}
\providecommand{\pandocbounded}[1]{#1}
\begin{document}
\begin{center}\Large\bfseries 交换对称成对损失：方法原型与完整实测\end{center}\section{研究定位}\label{ux7814ux7a76ux5b9aux4f4d}

本轮没有将模型改名当作原创，而是构造了一个固定预测器上的成对损失统计量，要求任意真假列交换时，被交换统计量反号、其他统计量不变。它是候选方法原型，不是已发表新算法，也尚未完成原创性系统检索。

研究已完成111个Olist开发样本实验、800份全新种子的模拟数据，每份都计算14种方法或消融，共1554个真实方法结果和11200个模拟方法结果。网络使用三个训练种子；未用旧模拟结果替代新计算。模型权重、日志、协议、数据摘要及失败检查可复现。

与此前不同，基础预测器不能用评分期响应早停。训练/早停完成后固定预测器，再在独立评分期产生排名。因此即使方法名称相同，本表也不能直接与旧实验未重新计算的RMSE混排。

\section{算法具体变化}\label{ux7b97ux6cd5ux5177ux4f53ux53d8ux5316}

\begin{enumerate}
\def\labelenumi{\arabic{enumi}.}
\tightlist
\item
  在每个真假数值对内取min/max，按固定随机掩码生成八组不区分真假的背景。
\item
  对特征j分别放入真实值和Knockoff值，计算平方预测损失差；其他变量始终使用相同的对称背景。
\item
  对背景、卖家内观测及卖家平均，形成W；标准化版再除以卖家均值标准误加常数。
\item
  使用既有Knockoff+阈值产生原生名单，同时单独报告Top-12/Top-8预算名单。
\end{enumerate}

本轮拆分普通替换、中点背景、随机对称背景、标准化四种评分，并分别接入同一批已训练XGBoost和官方TabM。另有SHAP、普通置换、Deep Lasso、评分集成对Lasso，以及使用同样训练样本的线性预测器。标准化不是t检验。

完整代数证明、先行文献以及独立样本/精确生成器前提见后附方法说明。符号翻转验证误差为零；这不证明Olist生成器精确。

\section{公平设计与解释边界}\label{ux516cux5e73ux8bbeux8ba1ux4e0eux89e3ux91caux8fb9ux754c}

所有真实方法共享同一卖家清单：30组互补50\%样本、20次70\%、20次80\%，另10组种子变化及完整参考。固定预算K=12，敏感性K=8。最终预测统一特征列顺序、模型参数、种子和时间窗口，相同变量集合复用相同预测。

真实面板共13,754行、2,723个卖家、31项指标；train/tune/rank/test分别为6,567/1,831/1,943/3,413行。响应为下一自然月签收确认商品金额的log1p，不含运费，不是按购买月回溯的最终送达GMV。RMSE、MAE和R²在log1p尺度计算，WAPE在金额尺度计算。静态卖家/商品信息的历史可见性仍有假设，不能将其包装为已经部署的实时预测系统。

新模拟800份分为八场景各100次，n=2000、p=31；训练1200行、评分400行、最终预测测试400行。七场景使用已知分布下的精确Knockoff，t错设场景故意使用不正确的高斯生成器。各场景和失败结果不按有利程度筛除。

成对Lasso只在评分样本拟合，其400行标签预算与基础模型在前1200行拟合的方式不同。所以主要还要看同样使用1200行训练、相同评分器的线性预测器对照。原生q=0.20，q=0.10敏感性在同一个W上重算。

真实数据存在估计Copula失配及同卖家跨期依赖，原生名单只称``名义阈值发现''，不得说实际FDR已被证明受控。2018年测试期曾多次查看，结果属探索性，不是独立外部确认。普通替换损失差不满足本轮要求的完整符号翻转，对它应用名义Knockoff阈值只是压力对照，不是CPI文献正式检验的复现。

\section{同预算总览}\label{ux540cux9884ux7b97ux603bux89c8}

{\def\LTcaptype{none} % do not increment counter
{
\footnotesize
\begin{longtable}[]{@{}
  >{\raggedright\arraybackslash}p{(\linewidth - 14\tabcolsep) * \real{0.240000}}
  >{\raggedright\arraybackslash}p{(\linewidth - 14\tabcolsep) * \real{0.108571}}
  >{\raggedright\arraybackslash}p{(\linewidth - 14\tabcolsep) * \real{0.108571}}
  >{\raggedright\arraybackslash}p{(\linewidth - 14\tabcolsep) * \real{0.108571}}
  >{\raggedright\arraybackslash}p{(\linewidth - 14\tabcolsep) * \real{0.108571}}
  >{\raggedright\arraybackslash}p{(\linewidth - 14\tabcolsep) * \real{0.108571}}
  >{\raggedright\arraybackslash}p{(\linewidth - 14\tabcolsep) * \real{0.108571}}
  >{\raggedright\arraybackslash}p{(\linewidth - 14\tabcolsep) * \real{0.108571}}@{}}
\toprule\noalign{}
\begin{minipage}[b]{\linewidth}\raggedright
方法
\end{minipage} & \begin{minipage}[b]{\linewidth}\raggedright
XGB-RMSE
\end{minipage} & \begin{minipage}[b]{\linewidth}\raggedright
TabM-RMSE
\end{minipage} & \begin{minipage}[b]{\linewidth}\raggedright
半样本J
\end{minipage} & \begin{minipage}[b]{\linewidth}\raggedright
Nogueira
\end{minipage} & \begin{minipage}[b]{\linewidth}\raggedright
FDP@12
\end{minipage} & \begin{minipage}[b]{\linewidth}\raggedright
Power@12
\end{minipage} & \begin{minipage}[b]{\linewidth}\raggedright
真实原生数
\end{minipage} \\
\midrule\noalign{}
\endhead
\bottomrule\noalign{}
\endlastfoot
评分集成对Lasso & 1.9488 & 1.9765 & 0.3569 & 0.2067 & 0.5203 & 0.6590 & 0 \\
线性预测器+标准化对称差 & 2.0784 & 2.0687 & 0.4046 & 0.2352 & 0.4860 & 0.7002 & 0 \\
XGBoost-SHAP & 1.8466 & 1.8571 & 0.7682 & 0.7439 & 0.2629 & 0.9678 & 无原生规则 \\
XGBoost-置换 & 1.8478 & 1.8612 & 0.7853 & 0.7666 & 0.2514 & 0.9817 & 无原生规则 \\
TabM-置换 & 1.8458 & 1.8586 & 0.7288 & 0.6841 & 0.2594 & 0.9720 & 无原生规则 \\
Deep Lasso & 1.8468 & 1.8567 & 0.8394 & 0.8012 & 0.2678 & 0.9620 & 无原生规则 \\
XGB-普通损失差 & 1.9286 & 1.9479 & 0.4237 & 0.2188 & 0.2669 & 0.9630 & 0 \\
XGB-中点对称差 & 1.9286 & 1.9479 & 0.4308 & 0.2269 & 0.2719 & 0.9570 & 0 \\
XGB-随机对称差 & 1.9286 & 1.9479 & 0.4300 & 0.2353 & 0.2697 & 0.9597 & 0 \\
XGB-标准化对称差 & 1.9323 & 1.9621 & 0.3387 & 0.1016 & 0.3171 & 0.9028 & 0 \\
TabM-普通损失差 & 2.0271 & 2.0294 & 0.3740 & 0.1947 & 0.2728 & 0.9560 & 0 \\
TabM-中点对称差 & 1.8577 & 1.8776 & 0.3934 & 0.2009 & 0.2769 & 0.9510 & 0 \\
TabM-随机对称差 & 1.8577 & 1.8776 & 0.3912 & 0.1993 & 0.2739 & 0.9547 & 0 \\
TabM-标准化对称差 & 2.0964 & 2.0886 & 0.3249 & 0.1248 & 0.3214 & 0.8977 & 0 \\
\end{longtable}
}
}

FDP@12和Power@12仅为六个精确生成器非零信号场景的等权平均；全无效和生成器错设另列。固定K时二者相关，不当作两份独立成功证据。真实原生发现为0也保留，预测采用常数均值作为空集基线。

\section{预定主候选与线性对照}\label{ux9884ux5b9aux4e3bux5019ux9009ux4e0eux7ebfux6027ux5bf9ux7167}

{\def\LTcaptype{none} % do not increment counter
{
\footnotesize
\begin{longtable}[]{@{}llllll@{}}
\toprule\noalign{}
场景 & TabM-FDP & TabM-Power & 线性-Power & Power差 & Holm p \\
\midrule\noalign{}
\endhead
\bottomrule\noalign{}
\endlastfoot
linear\_sparse & 0.1955 & 0.9700 & 1.0000 & -0.0300 & 0.666 \\
nonlinear\_strong & 0.1668 & 0.7940 & 0.0870 & +0.7070 & 1.11e-42 \\
nonlinear\_weak & 0.1636 & 0.4310 & 0.0750 & +0.3560 & 3.447e-15 \\
pure\_interactions & 0.1528 & 0.9340 & 0.0030 & +0.9310 & 5.535e-77 \\
high\_correlation & 0.1772 & 0.5840 & 0.8590 & -0.2750 & 4.367e-12 \\
mixed\_copula & 0.1506 & 0.8030 & 0.1750 & +0.6280 & 3.782e-34 \\
\end{longtable}
}
}

\textbf{已得到的判断。} 主候选在强非线性场景的Power为0.794，匹配线性对照为0.087，同时平均FDP为0.167；纯交互场景的Power为0.934对0.003。这支持利用非线性预测器发现线性评分遗漏的变量，但不证明优于所有已有非线性Knockoff方法。

\textbf{标准化改动没有获得支持。} 同一TabM预测器的未标准化随机对称差在强非线性场景Power为0.965，高于预定标准化版本。标准化版本在高相关线性场景为0.584，低于匹配线性对照的0.859。多个差异通过同场景Holm校正，不能把标准化描述成普遍提升检出率的创新。

\textbf{真实应用未形成优势。} 主实验Top-12、统一XGB评价下，主候选RMSE为2.0964，SHAP为1.8466；半样本J分别为0.3249和0.7682。主候选在这两个指标上均更差，不能以模拟Power的进步替代真实预测和稳定性的失败。

Holm校正在同场景、同指标、同规则的预定比较族内进行；不覆盖跨所有场景择优的结论。FDR是平均FDP的总体期望，单次FDP不必小于q。100次模拟只提供带Monte Carlo误差的估计，不是理论证明。

\section{全部原生发现结果}\label{ux5168ux90e8ux539fux751fux53d1ux73b0ux7ed3ux679c}

\subsection{global\_null}\label{global_null}

{\def\LTcaptype{none} % do not increment counter
{
\footnotesize
\begin{longtable}[]{@{}
  >{\raggedright\arraybackslash}p{(\linewidth - 12\tabcolsep) * \real{0.240000}}
  >{\raggedright\arraybackslash}p{(\linewidth - 12\tabcolsep) * \real{0.126667}}
  >{\raggedright\arraybackslash}p{(\linewidth - 12\tabcolsep) * \real{0.126667}}
  >{\raggedright\arraybackslash}p{(\linewidth - 12\tabcolsep) * \real{0.126667}}
  >{\raggedright\arraybackslash}p{(\linewidth - 12\tabcolsep) * \real{0.126667}}
  >{\raggedright\arraybackslash}p{(\linewidth - 12\tabcolsep) * \real{0.126667}}
  >{\raggedright\arraybackslash}p{(\linewidth - 12\tabcolsep) * \real{0.126667}}@{}}
\toprule\noalign{}
\begin{minipage}[b]{\linewidth}\raggedright
方法
\end{minipage} & \begin{minipage}[b]{\linewidth}\raggedright
FDP
\end{minipage} & \begin{minipage}[b]{\linewidth}\raggedright
MCSE
\end{minipage} & \begin{minipage}[b]{\linewidth}\raggedright
Power
\end{minipage} & \begin{minipage}[b]{\linewidth}\raggedright
发现数
\end{minipage} & \begin{minipage}[b]{\linewidth}\raggedright
非空比例
\end{minipage} & \begin{minipage}[b]{\linewidth}\raggedright
名单RMSE
\end{minipage} \\
\midrule\noalign{}
\endhead
\bottomrule\noalign{}
\endlastfoot
评分集成对Lasso & 0.0300 & 0.0171 & 不定义 & 0.18 & 0.03 & 0.9992 \\
线性预测器+标准化对称差 & 0.0100 & 0.0100 & 不定义 & 0.05 & 0.01 & 0.9989 \\
XGB-普通损失差 & 0.0600 & 0.0239 & 不定义 & 0.54 & 0.06 & 0.9996 \\
XGB-中点对称差 & 0.0600 & 0.0239 & 不定义 & 0.34 & 0.06 & 0.9998 \\
XGB-随机对称差 & 0.0400 & 0.0197 & 不定义 & 0.24 & 0.04 & 0.9992 \\
XGB-标准化对称差 & 0.0600 & 0.0239 & 不定义 & 0.43 & 0.06 & 0.9996 \\
TabM-普通损失差 & 0.0600 & 0.0239 & 不定义 & 0.36 & 0.06 & 0.9997 \\
TabM-中点对称差 & 0.0500 & 0.0219 & 不定义 & 0.32 & 0.05 & 0.9997 \\
TabM-随机对称差 & 0.0300 & 0.0171 & 不定义 & 0.24 & 0.03 & 0.9990 \\
TabM-标准化对称差 & 0.0500 & 0.0219 & 不定义 & 0.35 & 0.05 & 0.9998 \\
\end{longtable}
}
}

\subsection{linear\_sparse}\label{linear_sparse}

{\def\LTcaptype{none} % do not increment counter
{
\footnotesize
\begin{longtable}[]{@{}
  >{\raggedright\arraybackslash}p{(\linewidth - 12\tabcolsep) * \real{0.240000}}
  >{\raggedright\arraybackslash}p{(\linewidth - 12\tabcolsep) * \real{0.126667}}
  >{\raggedright\arraybackslash}p{(\linewidth - 12\tabcolsep) * \real{0.126667}}
  >{\raggedright\arraybackslash}p{(\linewidth - 12\tabcolsep) * \real{0.126667}}
  >{\raggedright\arraybackslash}p{(\linewidth - 12\tabcolsep) * \real{0.126667}}
  >{\raggedright\arraybackslash}p{(\linewidth - 12\tabcolsep) * \real{0.126667}}
  >{\raggedright\arraybackslash}p{(\linewidth - 12\tabcolsep) * \real{0.126667}}@{}}
\toprule\noalign{}
\begin{minipage}[b]{\linewidth}\raggedright
方法
\end{minipage} & \begin{minipage}[b]{\linewidth}\raggedright
FDP
\end{minipage} & \begin{minipage}[b]{\linewidth}\raggedright
MCSE
\end{minipage} & \begin{minipage}[b]{\linewidth}\raggedright
Power
\end{minipage} & \begin{minipage}[b]{\linewidth}\raggedright
发现数
\end{minipage} & \begin{minipage}[b]{\linewidth}\raggedright
非空比例
\end{minipage} & \begin{minipage}[b]{\linewidth}\raggedright
名单RMSE
\end{minipage} \\
\midrule\noalign{}
\endhead
\bottomrule\noalign{}
\endlastfoot
评分集成对Lasso & 0.1639 & 0.0213 & 1.0000 & 6.65 & 1.00 & 0.7438 \\
线性预测器+标准化对称差 & 0.1798 & 0.0182 & 1.0000 & 6.50 & 1.00 & 0.7438 \\
XGB-普通损失差 & 0.1845 & 0.0203 & 1.0000 & 6.72 & 1.00 & 0.7446 \\
XGB-中点对称差 & 0.1632 & 0.0196 & 1.0000 & 6.51 & 1.00 & 0.7442 \\
XGB-随机对称差 & 0.1865 & 0.0215 & 1.0000 & 6.80 & 1.00 & 0.7452 \\
XGB-标准化对称差 & 0.1903 & 0.0202 & 0.9980 & 6.78 & 1.00 & 0.7449 \\
TabM-普通损失差 & 0.2334 & 0.0229 & 1.0000 & 7.40 & 1.00 & 0.7450 \\
TabM-中点对称差 & 0.2013 & 0.0213 & 1.0000 & 6.94 & 1.00 & 0.7442 \\
TabM-随机对称差 & 0.1850 & 0.0219 & 1.0000 & 6.86 & 1.00 & 0.7447 \\
TabM-标准化对称差 & 0.1955 & 0.0205 & 0.9700 & 6.73 & 0.97 & 0.7584 \\
\end{longtable}
}
}

\subsection{nonlinear\_strong}\label{nonlinear_strong}

{\def\LTcaptype{none} % do not increment counter
{
\footnotesize
\begin{longtable}[]{@{}
  >{\raggedright\arraybackslash}p{(\linewidth - 12\tabcolsep) * \real{0.240000}}
  >{\raggedright\arraybackslash}p{(\linewidth - 12\tabcolsep) * \real{0.126667}}
  >{\raggedright\arraybackslash}p{(\linewidth - 12\tabcolsep) * \real{0.126667}}
  >{\raggedright\arraybackslash}p{(\linewidth - 12\tabcolsep) * \real{0.126667}}
  >{\raggedright\arraybackslash}p{(\linewidth - 12\tabcolsep) * \real{0.126667}}
  >{\raggedright\arraybackslash}p{(\linewidth - 12\tabcolsep) * \real{0.126667}}
  >{\raggedright\arraybackslash}p{(\linewidth - 12\tabcolsep) * \real{0.126667}}@{}}
\toprule\noalign{}
\begin{minipage}[b]{\linewidth}\raggedright
方法
\end{minipage} & \begin{minipage}[b]{\linewidth}\raggedright
FDP
\end{minipage} & \begin{minipage}[b]{\linewidth}\raggedright
MCSE
\end{minipage} & \begin{minipage}[b]{\linewidth}\raggedright
Power
\end{minipage} & \begin{minipage}[b]{\linewidth}\raggedright
发现数
\end{minipage} & \begin{minipage}[b]{\linewidth}\raggedright
非空比例
\end{minipage} & \begin{minipage}[b]{\linewidth}\raggedright
名单RMSE
\end{minipage} \\
\midrule\noalign{}
\endhead
\bottomrule\noalign{}
\endlastfoot
评分集成对Lasso & 0.0981 & 0.0182 & 0.1180 & 2.03 & 0.26 & 1.1743 \\
线性预测器+标准化对称差 & 0.0760 & 0.0169 & 0.0870 & 1.44 & 0.20 & 1.1842 \\
XGB-普通损失差 & 0.1946 & 0.0157 & 0.9980 & 12.97 & 1.00 & 0.8528 \\
XGB-中点对称差 & 0.1619 & 0.0135 & 0.9990 & 12.31 & 1.00 & 0.8533 \\
XGB-随机对称差 & 0.1879 & 0.0149 & 0.9990 & 12.78 & 1.00 & 0.8539 \\
XGB-标准化对称差 & 0.1701 & 0.0151 & 0.9240 & 11.72 & 1.00 & 0.8693 \\
TabM-普通损失差 & 0.1765 & 0.0159 & 0.9670 & 12.33 & 1.00 & 0.8585 \\
TabM-中点对称差 & 0.1808 & 0.0151 & 0.9570 & 12.20 & 1.00 & 0.8599 \\
TabM-随机对称差 & 0.1732 & 0.0155 & 0.9650 & 12.19 & 1.00 & 0.8589 \\
TabM-标准化对称差 & 0.1668 & 0.0166 & 0.7940 & 10.26 & 0.95 & 0.9021 \\
\end{longtable}
}
}

\subsection{nonlinear\_weak}\label{nonlinear_weak}

{\def\LTcaptype{none} % do not increment counter
{
\footnotesize
\begin{longtable}[]{@{}
  >{\raggedright\arraybackslash}p{(\linewidth - 12\tabcolsep) * \real{0.240000}}
  >{\raggedright\arraybackslash}p{(\linewidth - 12\tabcolsep) * \real{0.126667}}
  >{\raggedright\arraybackslash}p{(\linewidth - 12\tabcolsep) * \real{0.126667}}
  >{\raggedright\arraybackslash}p{(\linewidth - 12\tabcolsep) * \real{0.126667}}
  >{\raggedright\arraybackslash}p{(\linewidth - 12\tabcolsep) * \real{0.126667}}
  >{\raggedright\arraybackslash}p{(\linewidth - 12\tabcolsep) * \real{0.126667}}
  >{\raggedright\arraybackslash}p{(\linewidth - 12\tabcolsep) * \real{0.126667}}@{}}
\toprule\noalign{}
\begin{minipage}[b]{\linewidth}\raggedright
方法
\end{minipage} & \begin{minipage}[b]{\linewidth}\raggedright
FDP
\end{minipage} & \begin{minipage}[b]{\linewidth}\raggedright
MCSE
\end{minipage} & \begin{minipage}[b]{\linewidth}\raggedright
Power
\end{minipage} & \begin{minipage}[b]{\linewidth}\raggedright
发现数
\end{minipage} & \begin{minipage}[b]{\linewidth}\raggedright
非空比例
\end{minipage} & \begin{minipage}[b]{\linewidth}\raggedright
名单RMSE
\end{minipage} \\
\midrule\noalign{}
\endhead
\bottomrule\noalign{}
\endlastfoot
评分集成对Lasso & 0.0512 & 0.0171 & 0.0390 & 0.83 & 0.10 & 1.7252 \\
线性预测器+标准化对称差 & 0.0742 & 0.0177 & 0.0750 & 1.39 & 0.17 & 1.7212 \\
XGB-普通损失差 & 0.1784 & 0.0164 & 0.9160 & 11.81 & 1.00 & 1.5172 \\
XGB-中点对称差 & 0.1817 & 0.0156 & 0.9030 & 11.61 & 1.00 & 1.5204 \\
XGB-随机对称差 & 0.1733 & 0.0164 & 0.9020 & 11.63 & 0.99 & 1.5217 \\
XGB-标准化对称差 & 0.1926 & 0.0185 & 0.5540 & 8.04 & 0.77 & 1.5930 \\
TabM-普通损失差 & 0.1944 & 0.0185 & 0.7490 & 10.30 & 0.94 & 1.5428 \\
TabM-中点对称差 & 0.1823 & 0.0187 & 0.7580 & 10.34 & 0.94 & 1.5397 \\
TabM-随机对称差 & 0.1826 & 0.0174 & 0.7610 & 10.23 & 0.93 & 1.5411 \\
TabM-标准化对称差 & 0.1636 & 0.0189 & 0.4310 & 6.44 & 0.63 & 1.6157 \\
\end{longtable}
}
}

\subsection{pure\_interactions}\label{pure_interactions}

{\def\LTcaptype{none} % do not increment counter
{
\footnotesize
\begin{longtable}[]{@{}
  >{\raggedright\arraybackslash}p{(\linewidth - 12\tabcolsep) * \real{0.240000}}
  >{\raggedright\arraybackslash}p{(\linewidth - 12\tabcolsep) * \real{0.126667}}
  >{\raggedright\arraybackslash}p{(\linewidth - 12\tabcolsep) * \real{0.126667}}
  >{\raggedright\arraybackslash}p{(\linewidth - 12\tabcolsep) * \real{0.126667}}
  >{\raggedright\arraybackslash}p{(\linewidth - 12\tabcolsep) * \real{0.126667}}
  >{\raggedright\arraybackslash}p{(\linewidth - 12\tabcolsep) * \real{0.126667}}
  >{\raggedright\arraybackslash}p{(\linewidth - 12\tabcolsep) * \real{0.126667}}@{}}
\toprule\noalign{}
\begin{minipage}[b]{\linewidth}\raggedright
方法
\end{minipage} & \begin{minipage}[b]{\linewidth}\raggedright
FDP
\end{minipage} & \begin{minipage}[b]{\linewidth}\raggedright
MCSE
\end{minipage} & \begin{minipage}[b]{\linewidth}\raggedright
Power
\end{minipage} & \begin{minipage}[b]{\linewidth}\raggedright
发现数
\end{minipage} & \begin{minipage}[b]{\linewidth}\raggedright
非空比例
\end{minipage} & \begin{minipage}[b]{\linewidth}\raggedright
名单RMSE
\end{minipage} \\
\midrule\noalign{}
\endhead
\bottomrule\noalign{}
\endlastfoot
评分集成对Lasso & 0.0639 & 0.0221 & 0.0090 & 0.45 & 0.08 & 1.2239 \\
线性预测器+标准化对称差 & 0.0155 & 0.0109 & 0.0030 & 0.13 & 0.02 & 1.2233 \\
XGB-普通损失差 & 0.1793 & 0.0157 & 0.9830 & 12.57 & 1.00 & 1.0245 \\
XGB-中点对称差 & 0.1471 & 0.0128 & 0.9310 & 11.25 & 1.00 & 1.0306 \\
XGB-随机对称差 & 0.1690 & 0.0148 & 0.9370 & 11.75 & 1.00 & 1.0292 \\
XGB-标准化对称差 & 0.1620 & 0.0165 & 0.7090 & 9.34 & 0.88 & 1.0792 \\
TabM-普通损失差 & 0.1517 & 0.0145 & 1.0000 & 12.25 & 1.00 & 1.0125 \\
TabM-中点对称差 & 0.1565 & 0.0151 & 0.9980 & 12.31 & 1.00 & 1.0137 \\
TabM-随机对称差 & 0.1545 & 0.0144 & 0.9970 & 12.22 & 1.00 & 1.0157 \\
TabM-标准化对称差 & 0.1528 & 0.0149 & 0.9340 & 11.58 & 0.98 & 1.0279 \\
\end{longtable}
}
}

\subsection{high\_correlation}\label{high_correlation}

{\def\LTcaptype{none} % do not increment counter
{
\footnotesize
\begin{longtable}[]{@{}
  >{\raggedright\arraybackslash}p{(\linewidth - 12\tabcolsep) * \real{0.240000}}
  >{\raggedright\arraybackslash}p{(\linewidth - 12\tabcolsep) * \real{0.126667}}
  >{\raggedright\arraybackslash}p{(\linewidth - 12\tabcolsep) * \real{0.126667}}
  >{\raggedright\arraybackslash}p{(\linewidth - 12\tabcolsep) * \real{0.126667}}
  >{\raggedright\arraybackslash}p{(\linewidth - 12\tabcolsep) * \real{0.126667}}
  >{\raggedright\arraybackslash}p{(\linewidth - 12\tabcolsep) * \real{0.126667}}
  >{\raggedright\arraybackslash}p{(\linewidth - 12\tabcolsep) * \real{0.126667}}@{}}
\toprule\noalign{}
\begin{minipage}[b]{\linewidth}\raggedright
方法
\end{minipage} & \begin{minipage}[b]{\linewidth}\raggedright
FDP
\end{minipage} & \begin{minipage}[b]{\linewidth}\raggedright
MCSE
\end{minipage} & \begin{minipage}[b]{\linewidth}\raggedright
Power
\end{minipage} & \begin{minipage}[b]{\linewidth}\raggedright
发现数
\end{minipage} & \begin{minipage}[b]{\linewidth}\raggedright
非空比例
\end{minipage} & \begin{minipage}[b]{\linewidth}\raggedright
名单RMSE
\end{minipage} \\
\midrule\noalign{}
\endhead
\bottomrule\noalign{}
\endlastfoot
评分集成对Lasso & 0.1549 & 0.0156 & 0.7730 & 9.85 & 0.89 & 0.8352 \\
线性预测器+标准化对称差 & 0.1746 & 0.0115 & 0.8590 & 10.75 & 0.95 & 0.8088 \\
XGB-普通损失差 & 0.1851 & 0.0170 & 0.7280 & 9.82 & 0.89 & 0.8412 \\
XGB-中点对称差 & 0.1712 & 0.0163 & 0.7020 & 9.34 & 0.85 & 0.8613 \\
XGB-随机对称差 & 0.1779 & 0.0171 & 0.7290 & 9.77 & 0.89 & 0.8418 \\
XGB-标准化对称差 & 0.1514 & 0.0204 & 0.3700 & 5.65 & 0.53 & 1.0223 \\
TabM-普通损失差 & 0.1662 & 0.0165 & 0.8480 & 11.01 & 0.95 & 0.8086 \\
TabM-中点对称差 & 0.1617 & 0.0163 & 0.8440 & 10.87 & 0.96 & 0.8059 \\
TabM-随机对称差 & 0.1686 & 0.0163 & 0.8510 & 11.04 & 0.97 & 0.8000 \\
TabM-标准化对称差 & 0.1772 & 0.0197 & 0.5840 & 8.46 & 0.73 & 0.9239 \\
\end{longtable}
}
}

\subsection{mixed\_copula}\label{mixed_copula}

{\def\LTcaptype{none} % do not increment counter
{
\footnotesize
\begin{longtable}[]{@{}
  >{\raggedright\arraybackslash}p{(\linewidth - 12\tabcolsep) * \real{0.240000}}
  >{\raggedright\arraybackslash}p{(\linewidth - 12\tabcolsep) * \real{0.126667}}
  >{\raggedright\arraybackslash}p{(\linewidth - 12\tabcolsep) * \real{0.126667}}
  >{\raggedright\arraybackslash}p{(\linewidth - 12\tabcolsep) * \real{0.126667}}
  >{\raggedright\arraybackslash}p{(\linewidth - 12\tabcolsep) * \real{0.126667}}
  >{\raggedright\arraybackslash}p{(\linewidth - 12\tabcolsep) * \real{0.126667}}
  >{\raggedright\arraybackslash}p{(\linewidth - 12\tabcolsep) * \real{0.126667}}@{}}
\toprule\noalign{}
\begin{minipage}[b]{\linewidth}\raggedright
方法
\end{minipage} & \begin{minipage}[b]{\linewidth}\raggedright
FDP
\end{minipage} & \begin{minipage}[b]{\linewidth}\raggedright
MCSE
\end{minipage} & \begin{minipage}[b]{\linewidth}\raggedright
Power
\end{minipage} & \begin{minipage}[b]{\linewidth}\raggedright
发现数
\end{minipage} & \begin{minipage}[b]{\linewidth}\raggedright
非空比例
\end{minipage} & \begin{minipage}[b]{\linewidth}\raggedright
名单RMSE
\end{minipage} \\
\midrule\noalign{}
\endhead
\bottomrule\noalign{}
\endlastfoot
评分集成对Lasso & 0.1177 & 0.0211 & 0.1640 & 2.78 & 0.33 & 1.1642 \\
线性预测器+标准化对称差 & 0.0903 & 0.0172 & 0.1750 & 2.56 & 0.33 & 1.1590 \\
XGB-普通损失差 & 0.1680 & 0.0147 & 0.9980 & 12.48 & 1.00 & 0.8517 \\
XGB-中点对称差 & 0.1750 & 0.0154 & 0.9940 & 12.56 & 1.00 & 0.8523 \\
XGB-随机对称差 & 0.1849 & 0.0161 & 0.9940 & 12.80 & 1.00 & 0.8526 \\
XGB-标准化对称差 & 0.1684 & 0.0169 & 0.8940 & 11.55 & 0.97 & 0.8799 \\
TabM-普通损失差 & 0.1888 & 0.0162 & 0.9810 & 12.69 & 1.00 & 0.8546 \\
TabM-中点对称差 & 0.1612 & 0.0163 & 0.9640 & 12.10 & 1.00 & 0.8572 \\
TabM-随机对称差 & 0.1518 & 0.0144 & 0.9630 & 11.78 & 1.00 & 0.8573 \\
TabM-标准化对称差 & 0.1506 & 0.0146 & 0.8030 & 10.00 & 0.97 & 0.8999 \\
\end{longtable}
}
}

\subsection{misspecified\_t}\label{misspecified_t}

{\def\LTcaptype{none} % do not increment counter
{
\footnotesize
\begin{longtable}[]{@{}
  >{\raggedright\arraybackslash}p{(\linewidth - 12\tabcolsep) * \real{0.240000}}
  >{\raggedright\arraybackslash}p{(\linewidth - 12\tabcolsep) * \real{0.126667}}
  >{\raggedright\arraybackslash}p{(\linewidth - 12\tabcolsep) * \real{0.126667}}
  >{\raggedright\arraybackslash}p{(\linewidth - 12\tabcolsep) * \real{0.126667}}
  >{\raggedright\arraybackslash}p{(\linewidth - 12\tabcolsep) * \real{0.126667}}
  >{\raggedright\arraybackslash}p{(\linewidth - 12\tabcolsep) * \real{0.126667}}
  >{\raggedright\arraybackslash}p{(\linewidth - 12\tabcolsep) * \real{0.126667}}@{}}
\toprule\noalign{}
\begin{minipage}[b]{\linewidth}\raggedright
方法
\end{minipage} & \begin{minipage}[b]{\linewidth}\raggedright
FDP
\end{minipage} & \begin{minipage}[b]{\linewidth}\raggedright
MCSE
\end{minipage} & \begin{minipage}[b]{\linewidth}\raggedright
Power
\end{minipage} & \begin{minipage}[b]{\linewidth}\raggedright
发现数
\end{minipage} & \begin{minipage}[b]{\linewidth}\raggedright
非空比例
\end{minipage} & \begin{minipage}[b]{\linewidth}\raggedright
名单RMSE
\end{minipage} \\
\midrule\noalign{}
\endhead
\bottomrule\noalign{}
\endlastfoot
评分集成对Lasso & 0.1546 & 0.0286 & 0.0790 & 2.12 & 0.24 & 1.2047 \\
线性预测器+标准化对称差 & 0.0755 & 0.0208 & 0.0370 & 1.00 & 0.13 & 1.2168 \\
XGB-普通损失差 & 0.0968 & 0.0152 & 0.5060 & 6.36 & 0.69 & 1.0530 \\
XGB-中点对称差 & 0.1382 & 0.0167 & 0.5090 & 6.80 & 0.70 & 1.0509 \\
XGB-随机对称差 & 0.1163 & 0.0154 & 0.4990 & 6.42 & 0.68 & 1.0497 \\
XGB-标准化对称差 & 0.1419 & 0.0189 & 0.3660 & 5.35 & 0.57 & 1.0911 \\
TabM-普通损失差 & 0.1574 & 0.0198 & 0.4060 & 5.63 & 0.69 & 1.0560 \\
TabM-中点对称差 & 0.1390 & 0.0192 & 0.3710 & 5.02 & 0.67 & 1.0655 \\
TabM-随机对称差 & 0.1551 & 0.0202 & 0.3970 & 5.52 & 0.69 & 1.0580 \\
TabM-标准化对称差 & 0.1783 & 0.0228 & 0.2740 & 4.49 & 0.53 & 1.1149 \\
\end{longtable}
}
}

全无效场景Power没有定义；其FDR等于选出非空集合的概率。准确二项区间在simulation\_summary.csv中报告，不把100次全未发现说成真实错误率必为0。发现数不同的原生名单与强制12项名单不混比。

主候选全无效场景为100次中5次非空，估计FDR=0.050，准确95\%区间为{[}0.016, 0.113{]}。线性稀疏场景的TabM中点平均FDP略超过0.20，仍须连同Monte Carlo标准误如实保留，不能宣称每个有限模拟均值都严格低于目标。

\subsection{同一主候选的阈值敏感性}\label{ux540cux4e00ux4e3bux5019ux9009ux7684ux9608ux503cux654fux611fux6027}

{\def\LTcaptype{none} % do not increment counter
{
\footnotesize
\begin{longtable}[]{@{}llllll@{}}
\toprule\noalign{}
场景 & q=.2 FDP & q=.2 Power & q=.1 FDP & q=.1 Power & q=.1发现数 \\
\midrule\noalign{}
\endhead
\bottomrule\noalign{}
\endlastfoot
global\_null & 0.0500 & 不定义 & 0.0100 & 不定义 & 0.11 \\
linear\_sparse & 0.1955 & 0.9700 & 0.0167 & 0.0300 & 0.34 \\
nonlinear\_strong & 0.1668 & 0.7940 & 0.0607 & 0.2740 & 3.45 \\
nonlinear\_weak & 0.1636 & 0.4310 & 0.0399 & 0.0900 & 1.37 \\
pure\_interactions & 0.1528 & 0.9340 & 0.0588 & 0.5260 & 5.96 \\
high\_correlation & 0.1772 & 0.5840 & 0.0477 & 0.1880 & 2.54 \\
mixed\_copula & 0.1506 & 0.8030 & 0.0602 & 0.3270 & 3.96 \\
misspecified\_t & 0.1783 & 0.2740 & 0.0325 & 0.0440 & 0.82 \\
\end{longtable}
}
}

这些是相同W的不同阈值，不是重新训练的模型。q=.1至少需要10个正发现才可能非空，稀疏信号场景可能因此损失Power。t错设行只有经验意义，即使平均FDP低于目标也不能恢复不存在的理论保证。

\section{真实生成器诊断}\label{ux771fux5b9eux751fux6210ux5668ux8bcaux65ad}

{\def\LTcaptype{none} % do not increment counter
{
\footnotesize
\begin{longtable}[]{@{}ll@{}}
\toprule\noalign{}
交换比例 & 卖家隔离AUC \\
\midrule\noalign{}
\endhead
\bottomrule\noalign{}
\endlastfoot
0.25 & 0.9514 \\
0.5 & 0.6470 \\
1.0 & 0.8998 \\
\end{longtable}
}
}

分类训练与测试按卖家分离，原样本及交换样本成对划入同一侧。该AUC仅为工程诊断，不提供检验p值；接近0.5也不是交换性证明。因为评分所用卖家可能出现在模型训练历史中，独立样本理论不能直接外推。

\section{怎样评价论文可行性}\label{ux600eux6837ux8bc4ux4ef7ux8bbaux6587ux53efux884cux6027}

有明确研究问题、算法定义、代数性质、机制消融和完整数据实测，比不断拼接网络更具可讨论性。但是否足以成为学位论文方法贡献，需要分别判断：原创性是否与先行文献区分、目标场景是否真实改善、真实数据是否有可使用结果。不能用大批训练次数替代任何一项。

本轮的明确判断是：可以作为``模型无关损失评分的对称化、非线性检出及生成器敏感性''方法研究原型，但尚不能作为``已证明原创且在电商数据上全面改进''的定稿成果。预定标准化版本不宜成为成功创新的核心叙事；简单对称化是否与CPI、Semi-knockoffs等文献实质不同，还需逐公式对照，并补足近期非线性受控筛选基线和第二个未使用过的真实数据集。

可以讨论的暂定题目是《面向电商经营指标筛选的交换对称损失统计量及其稳健性研究》。这个题目描述研究问题，不承诺方法优于一切基线。当前不能保证其创新程度满足学位要求，更不能保证答辩通过。

\subsection{计算成本与复现审计}\label{ux8ba1ux7b97ux6210ux672cux4e0eux590dux73b0ux5ba1ux8ba1}

{\def\LTcaptype{none} % do not increment counter
{
\footnotesize
\begin{longtable}[]{@{}llll@{}}
\toprule\noalign{}
运行范围 & 完成案例 & 每案例秒数中位数 & 案例耗时合计/小时 \\
\midrule\noalign{}
\endhead
\bottomrule\noalign{}
\endlastfoot
real & 111 & 48.2 & 1.75 \\
simulations & 800 & 15.4 & 3.27 \\
real\_mvr & 111 & 57.9 & 1.88 \\
\end{longtable}
}
}

每案例同时训练和计算多种共享预测器/统计量。案例耗时合计是各进程计时之和，不是整轮墙钟耗时或CPU核心小时，不能拆算为某方法独有成本。本轮未保存逐评分方法的独立计时，因此不提供虚构的速度排名。

审计对预测RMSE、MAE、R²、WAPE从逐行预测重算，并检查样本哈希、两种原生阈值、有限分数和相同集合预测一致性。50\%稳定性区间按30组互补样本计算，不把60个半样本当独立重复。完整补充指标、配对区间及代码摘要保存在results目录；生成报告没有重调模型参数或改动冻结协议。

逐行预测最初含float32输出，CSV重读后按float64复算金额指数函数会产生约2e-8的差异；核对采用rtol=2e-7、atol=1e-9，不影响报告四位小数。NumPy/BLAS矩阵乘法警告保留在运行记录中，不通过关闭警告隐藏。附加的10万行精确高斯检查同时核验随机背景协方差，并用不调用BLAS的显式einsum交叉核对采样矩阵乘法；它不是真实Olist生成器合格证明。

\section{事后追加：MVR生成器敏感性}\label{ux4e8bux540eux8ffdux52a0mvrux751fux6210ux5668ux654fux611fux6027}

主协议的真实等相关S约0.00334，约49.4\%的真假数值完全相同，比较信息很弱。一次参考样本检查已看到MVR的名义发现数及测试RMSE变化，随后才制定mvr\_extension\_protocol.json，追加相同111个真实案例及两种评价器的Top-12分析，不改动前述主表。追加动机受到已见结果启发；不能把``未搜索生成器参数或修改q''误写成``生成器选择完全未受测试表现影响''。

{\def\LTcaptype{none} % do not increment counter
{
\footnotesize
\begin{longtable}[]{@{}
  >{\raggedright\arraybackslash}p{(\linewidth - 12\tabcolsep) * \real{0.240000}}
  >{\raggedright\arraybackslash}p{(\linewidth - 12\tabcolsep) * \real{0.126667}}
  >{\raggedright\arraybackslash}p{(\linewidth - 12\tabcolsep) * \real{0.126667}}
  >{\raggedright\arraybackslash}p{(\linewidth - 12\tabcolsep) * \real{0.126667}}
  >{\raggedright\arraybackslash}p{(\linewidth - 12\tabcolsep) * \real{0.126667}}
  >{\raggedright\arraybackslash}p{(\linewidth - 12\tabcolsep) * \real{0.126667}}
  >{\raggedright\arraybackslash}p{(\linewidth - 12\tabcolsep) * \real{0.126667}}@{}}
\toprule\noalign{}
\begin{minipage}[b]{\linewidth}\raggedright
方法
\end{minipage} & \begin{minipage}[b]{\linewidth}\raggedright
XGB-RMSE
\end{minipage} & \begin{minipage}[b]{\linewidth}\raggedright
TabM-RMSE
\end{minipage} & \begin{minipage}[b]{\linewidth}\raggedright
半样本J
\end{minipage} & \begin{minipage}[b]{\linewidth}\raggedright
Nogueira
\end{minipage} & \begin{minipage}[b]{\linewidth}\raggedright
名义发现数
\end{minipage} & \begin{minipage}[b]{\linewidth}\raggedright
半样本空集率
\end{minipage} \\
\midrule\noalign{}
\endhead
\bottomrule\noalign{}
\endlastfoot
评分集成对Lasso & 1.8941 & 1.8790 & 0.5216 & 0.3650 & 0 & 0.317 \\
线性预测器+标准化对称差 & 1.8825 & 1.8849 & 0.6017 & 0.4680 & 10 & 0.367 \\
XGBoost-SHAP & 1.8466 & 1.8571 & 0.7682 & 0.7439 & 无 & 不适用 \\
XGBoost-置换 & 1.8478 & 1.8612 & 0.7853 & 0.7666 & 无 & 不适用 \\
TabM-置换 & 1.8458 & 1.8586 & 0.7288 & 0.6841 & 无 & 不适用 \\
Deep Lasso & 1.8468 & 1.8567 & 0.8394 & 0.8012 & 无 & 不适用 \\
XGB-普通损失差 & 1.8484 & 1.8577 & 0.5515 & 0.4225 & 6 & 0.067 \\
XGB-中点对称差 & 1.8662 & 1.8727 & 0.4982 & 0.4660 & 5 & 0.333 \\
XGB-随机对称差 & 1.8521 & 1.8668 & 0.5119 & 0.4541 & 6 & 0.117 \\
XGB-标准化对称差 & 1.8479 & 1.8680 & 0.4817 & 0.3451 & 10 & 0.250 \\
TabM-普通损失差 & 1.8899 & 1.8974 & 0.4875 & 0.4268 & 5 & 0.167 \\
TabM-中点对称差 & 1.8431 & 1.8575 & 0.4988 & 0.4509 & 6 & 0.100 \\
TabM-随机对称差 & 1.8630 & 1.8732 & 0.4874 & 0.4675 & 6 & 0.100 \\
TabM-标准化对称差 & 1.8493 & 1.8645 & 0.4696 & 0.3947 & 9 & 0.083 \\
\end{longtable}
}
}

{\def\LTcaptype{none} % do not increment counter
{
\footnotesize
\begin{longtable}[]{@{}ll@{}}
\toprule\noalign{}
交换比例 & 卖家隔离AUC \\
\midrule\noalign{}
\endhead
\bottomrule\noalign{}
\endlastfoot
0.25 & 0.8782 \\
0.5 & 0.9139 \\
1.0 & 0.9467 \\
\end{longtable}
}
}

这是观察主实验结果后追加的诊断，不是预先注册确认。MVR为既有方法，不列为本轮原创；相同样本、基础模型、阈值保持不变，四项不使用Knockoff的基线排名逐案例与主实验一致。只重新检验真实样本，不把前述精确等相关模拟结果误写成MVR实测FDR。发现数增加不证明原生FDR正确；估计Copula及跨期卖家依赖的限制依然存在。

主候选在MVR下的XGB评价RMSE为1.8493，半样本J为0.4696；TabM中点版本分别为1.8431和0.4988。中点版本只属于已登记的机制消融，不能事后改称预定主候选。主实验、追加实验不能各抽一个最好的指标拼成``总体最优方法''。

为避免把小幅数值差当成确定收益，另对追加实验全部六个普通/中点/随机背景版本补算相对SHAP的配对区间，标记为事后诊断，不改变主比较族。中点TabM相对SHAP的XGB评价RMSE差为-0.0035，卖家簇Bootstrap点态95\%区间为{[}-0.0118, +0.0053{]}。这不是新测试集上的确认性检验；区间包含0时不能声称已证实优于SHAP，也不能声称两者已证明等价。

\section{附录：方法定义与证明边界}\label{ux9644ux5f55ux65b9ux6cd5ux5b9aux4e49ux4e0eux8bc1ux660eux8fb9ux754c}

\subsection{研究问题}\label{ux7814ux7a76ux95eeux9898}

在训练好的预测模型之外，能否构造一个不依赖网络权重解释、直接比较预测损失、同时满足Knockoff完整符号翻转性质的变量筛选统计量？

这是本轮实际研究的问题，不是``找一个更大的网络打败XGBoost''。候选方法简称SCPL（swap-canonical paired loss）。名称只是工程标识，不构成原创性证明。本轮独立编写的部分是对称背景、成对损失、固定模型评分及标准化之间的具体组合与验证；Knockoff、条件预测重要性、反对称统计量和集成网络本身均来自既有工作。

\subsection{与既有方法的关系}\label{ux4e0eux65e2ux6709ux65b9ux6cd5ux7684ux5173ux7cfb}

{\def\LTcaptype{none} % do not increment counter
{
\footnotesize
\begin{longtable}[]{@{}
  >{\raggedright\arraybackslash}p{(\linewidth - 4\tabcolsep) * \real{0.3333}}
  >{\raggedright\arraybackslash}p{(\linewidth - 4\tabcolsep) * \real{0.3333}}
  >{\raggedright\arraybackslash}p{(\linewidth - 4\tabcolsep) * \real{0.3333}}@{}}
\toprule\noalign{}
\begin{minipage}[b]{\linewidth}\raggedright
真实文献
\end{minipage} & \begin{minipage}[b]{\linewidth}\raggedright
核心内容
\end{minipage} & \begin{minipage}[b]{\linewidth}\raggedright
本轮不能冒称为原创的部分
\end{minipage} \\
\midrule\noalign{}
\endhead
\bottomrule\noalign{}
\endlastfoot
Candès等，Panning for gold: Model-X knockoffs for high dimensional controlled variable selection，JRSS B，2018，DOI 10.1111/rssb.12265 & Model-X条件独立、交换性、Knockoff+ & 错误率控制框架 \\
Lu等，DeepPINK: reproducible feature selection in deep neural networks，NeurIPS 2018，arXiv:1809.01185 & 成对输入与深度选择 & 真假特征进入同一深度模型 \\
Watson和Wright，Testing Conditional Independence in Supervised Learning Algorithms，Machine Learning 110:2107--2129，2021，DOI 10.1007/s10994-021-06030-6 & 条件预测影响CPI与损失差 & 用Knockoff替换输入后比较预测损失 \\
Oh和Lee，DeepPIG，Scientific Reports 14:15582，2024，DOI 10.1038/s41598-024-66061-6 & 成对层、随机门与Knockoff & 成对输入加门控这一方向 \\
Gorishniy等，TabM: Advancing Tabular Deep Learning with Parameter-Efficient Ensembling，ICLR 2025，arXiv:2410.24210 & 参数高效表格集成 & TabM预测架构 \\
Zou和Tian，GRIP2: A Robust and Powerful Deep Knockoff Method for Feature Selection，arXiv:2602.00218，2026预印本 & 对正则化曲面上的特征活动聚合 & 深度统计量和轨迹聚合已有相邻研究 \\
Reyero-Lobo等，Semi-knockoffs: a model-agnostic conditional independence testing method with finite-sample guarantees，ICML 2026，arXiv:2601.23124v2 & 对预训练模型的条件独立检验 & 模型无关的受控筛选已有近期研究 \\
\end{longtable}
}
}

上述来源已通过期刊/会议页或arXiv作者摘要核验。这里只进行了针对性先行研究排查，没有完成足以证明``此前无人提出''的系统检索。本文没有冒充复现DeepPIG、GRIP2或Semi-knockoffs；它们是必须讨论的相邻文献，而非已完成实证基线。

\subsection{数据分离}\label{ux6570ux636eux5206ux79bb}

先用训练数据A拟合预测器f，内部早停后在A内固定轮次重拟合。随后在不参与训练、早停或调参的评分数据B上生成Knockoff并计算统计量。最终测试数据C只计算名单的下游预测表现。

在模拟中，A有1200行，其中前1000行训练、后200行选轮次，B有400行，C有400行。Olist沿用train/tune/rank/test时间切分，rank期只评分，不参与基础预测器的早停。

Olist同一卖家跨期相关，所以A与B不满足理论中的独立样本假设。真实数据分析只能是经验探索；用卖家抽样不能消除这一理论缺口。

\subsection{对称背景与损失对比}\label{ux5bf9ux79f0ux80ccux666fux4e0eux635fux5931ux5bf9ux6bd4}

对B中每个观测i和特征j，有真实值x与Knockoff值x̃。先取：

\begin{verbatim}
低值 Lij = min(xij, x̃ij)
高值 Uij = max(xij, x̃ij)
\end{verbatim}

生成四组与响应无关的随机0/1掩码及其补掩码，共八个背景C。每个背景对每个变量从低值、高值中选一个。这个背景只依赖无序数值对，不依赖哪一个被命名为``真实''。

评价特征j时，其他变量使用背景值，分别把第j列设为真实值与Knockoff值：

\begin{verbatim}
dij = 八个背景下的平均 [
  使用Knockoff第j列的平方误差
  - 使用真实第j列的平方误差
]
\end{verbatim}

正值表示：在这些对称背景中，真实第j列比其Knockoff对应值更有利于预测。

首先在每个卖家内平均d，再对卖家平均形成W。标准化版以W除以卖家均值的标准误加固定正数。模拟每行是独立组，因此退化为普通行平均。

``标准误''在此仅用于无量纲排序，不进行Student t检验，不宣称统计量服从t分布。用卖家聚合不会自动赋予面板FDR保证。

\subsection{符号翻转的代数命题}\label{ux7b26ux53f7ux7ffbux8f6cux7684ux4ee3ux6570ux547dux9898}

\textbf{命题（固定预测器下的路径级符号翻转）。} 假定f、预处理、分组、掩码和任何超参数在交换评分数据的真假变量之前已固定。对任意变量集合S，交换每个j属于S的真实列与对应Knockoff列，则：

\[
W_j([X,\widetilde X]_{\mathrm{swap}(S)},Y)
=\begin{cases}
-W_j([X,\widetilde X],Y),&j\in S,\\
W_j([X,\widetilde X],Y),&j\notin S.
\end{cases}
\]

\textbf{证明。} 交换不改变任意行任意列的min/max，因此所有背景不变。对于j不在S中，比较时自身真假值未变，其他列仍使用相同背景，两次预测及其差均不变。对于j在S中，自身两次预测的输入相互交换，因此每行、每背景的损失差反号。卖家内及卖家间平均保留反号。标准化版的组间标准差在整体反号时不变，所以比值也只反号。证毕。

这一证明不要求f正确拟合真实响应函数。它证明的是统计量的对称性，而不是预测精度、检出率、模型收敛或新生成器正确性。min/max背景也可替换为其他对交换不变的背景；本轮另跑中点背景作为消融。

\textbf{补充性质（精确生成器下随机背景的边际分布）。} 对独立同分布观测，若每行真假变量满足精确逐坐标交换性，且对掩码随机性取平均，则任一随机min/max背景的整行边际分布与原X相同。原因是给定真假数值，每坐标以各一半概率取min/max，等价于以各一半概率取真实值/Knockoff值。因此背景分布是全部坐标交换子集的等权混合；每个交换子集下的原侧向量都与X同分布，其混合也同分布。补掩码不改变各背景的边际分布，但背景之间不独立。

这一性质需要对掩码随机性平均，不能说任一固定min/max取值都来自原联合分布，也不说明有限八背景没有Monte Carlo误差。它只涉及输入边际，不意味着背景与Y的关系不变，不证明Power最优。中点背景通常没有这个性质。Olist的估计Copula失配时同样不能套用。该性质是对现有实现的理论说明，不是看结果后改变算法，也不主张这一群对称化思想的优先权。

\subsection{可以调用的既有FDR结论}\label{ux53efux4ee5ux8c03ux7528ux7684ux65e2ux6709fdrux7ed3ux8bba}

还必须另外假设：

\begin{enumerate}
\def\labelenumi{\arabic{enumi}.}
\tightlist
\item
  条件于A，B中的样本来自目标分布并满足所用Model-X框架所需的样本条件；
\item
  B的Knockoff与响应在给定X及A时条件独立；
\item
  B的联合真假分布满足所需的精确交换性；
\item
  f及评分超参数不使用B的响应训练或选优；
\item
  掩码独立产生，统计量满足上面的完整符号翻转。
\end{enumerate}

在这些条件下，可调用既有Knockoff+结果，在目标水平q控制FDR。这里不是新提出一个FDR定理，而是把本轮统计量代入既有定理的一项充分条件检查。

反对称性不能补救错误的Knockoff生成器。在Olist的估计Copula和跨期卖家依赖中，上述假设未被证明。本轮故意加入t分布数据配高斯生成器，检验错设时可能出现的风险，不会将其平均FDP也标为理论保证。

\subsection{普通替换为何不能直接套同一结论}\label{ux666eux901aux66ffux6362ux4e3aux4f55ux4e0dux80fdux76f4ux63a5ux5957ux540cux4e00ux7ed3ux8bba}

普通损失差把其他列固定为``真实列''。交换某一个变量后，该列也进入其他特征的比较背景，其他W可能随之改变。因此，``自己那一列的差反号''不够，还要保证其他统计量不变。

本轮普通替换版本仅作为数值和预测比较的对照。若对它使用名义Knockoff+阈值，结果必须明确标记``假设未获保证的压力测试''，不是论文CPI原检验的完整复现，也不能据此否定CPI文献中的其他检验程序。

\subsection{代价、弱点与预算}\label{ux4ee3ux4ef7ux5f31ux70b9ux4e0eux9884ux7b97}

\begin{itemize}
\tightlist
\item
  中点背景可能落在训练分布外；随机背景在精确生成器下具有上面的边际同分布性质，但改变了与响应的关系，有限背景近似也有误差。生成器错设时随机背景仍可能失配；符号翻转成立不意味着检出率更高。
\item
  原生成器准确性仍是难点；没有创造``任意分布都有效''的假变量。
\item
  需要独立评分样本，牺牲可用于模型训练的数据量。
\item
  推断计算约需2×特征数×背景数次批预测；本轮八背景为固定预算，不因结果好坏改动。
\item
  成对Lasso只在B中拟合，和利用A训练的模型信息用法不同。额外的线性预测器使用同样A并用同样评分规则，避免把A的额外信息单独归给非线性改动。
\item
  本轮原生q=0.20，同时报告q=0.10。Knockoff+至少需要ceil(1/q)个发现才可能非空：不能把放宽q带来的Power变化写成算法创新。
\item
  Top-K只是预算名单，不继承原生FDR保证。全无效变量场景没有Power，必须记缺失而非人为填零。
\end{itemize}

\subsection{成为论文方法贡献之前}\label{ux6210ux4e3aux8bbaux6587ux65b9ux6cd5ux8d21ux732eux4e4bux524d}

本轮比单纯拼接更明确的一点，是有一个可检查的统计量设计问题、公式、代数性质、对称性反例和拆分消融。但这不自动等于原创性已经成立。仍需要更完整的相关方法比较、生成器稳健性分析，以及新的真实数据验证。

用户需要的是可辩护的研究，而不是``自称独创''。如果实际结果不支持Power、预测或稳定性改进，就按失败或有边界的结果报告。

\end{document}
